\ifdefined\gkver\else\def\gkver{student}\fi
\documentclass[\gkver]{gaokaozhenti}
\begin{document}

\chapter{2025年贵州}

\begin{enumerate}
\item
%% number: 1
%% paperName: 2025年贵州高考第 1 题 4 分
%% typeId: 1
%% type: 单选题
%% chapter: 直线运动
%% point: 运动图象
%% method: 无
%% score: 4
%% degree: 850
%% duplicateId: 0
%% body:
为测试人形机器人的稳定性和灵活性，让人形机器人按指令在一条直线上“跑步”，如图~\ref{2025贵州01a}所示。人形机器人在一段时间内的位置-时间图像如图~\ref{2025贵州01b}所示，设该机器人在 $0 \sim 10 \Us$ 内的位移为 $\Delta x_{1}$、速度为 $v_{1}$，$20 \sim 30 \Us$ 内的位移为 $\Delta x_{2}$、速度为 $v_{2}$，则\xzanswer{D}

\twopicture[labelA=2025贵州01a, labelB=2025贵州01b, wA=2.6cm, wB=4.6cm]{figs/01a.png}{figs/01b.png}

\fourchoices[answer=D]
{$\Delta x_{1} < \Delta x_{2}$，$v_{1} < v_{2}$}
{$\Delta x_{1} < \Delta x_{2}$，$v_{1} > v_{2}$}
{$\Delta x_{1} > \Delta x_{2}$，$v_{1} < v_{2}$}
{$\Delta x_{1} > \Delta x_{2}$，$v_{1} > v_{2}$}

%% answer: D
%% memo:
\memoanswer{
根据图像可知，$0 \sim 10 \Us$ 内的位移为 $\Delta x_{1} = 20 \Um$

速度为 $v_{1} = \frac{\Delta x_{1}}{\Delta t_{1}} = 2 \Ums$

$20 \sim 30 \Us$ 内的位移为 $\Delta x_{2} = 30 \Um - 20 \Um = 10 \Um$

速度为 $v_{2} = \frac{\Delta x_{2}}{\Delta t_{2}} = 1 \Ums$

所以 $\Delta x_{1} > \Delta x_{2}$，$v_{1} > v_{2}$

故选 D。
}

\item
%% number: 2
%% paperName: 2025年贵州高考第 2 题 4 分
%% typeId: 1
%% type: 单选题
%% chapter: 运动和力
%% point: 动量守恒定律
%% method: 无
%% score: 4
%% degree: 750
%% duplicateId: 0
%% body:
甲、乙两运动员在光滑水平冰面上进行滑冰训练。以速度 $v_{0}$ 运动的甲推一下静止在其正前方的乙，刚分开时，甲、乙的运动方向与 $v_{0}$ 的方向相同，且甲的速度为 $\frac{v_{0}}{3}$，则刚分开时，甲、乙的动量大小之比为\xzanswer{C}

\fourchoices[answer=C]
{3:1}
{2:1}
{1:2}
{1:3}

%% answer: C
%% memo:
\memoanswer{
根据 $p = mv$，设甲的初始动量为 $p_{0}$，则刚分开时，甲的末动量为 $\frac{1}{3}p_{0}$，甲的动量变化量大小为 $\Delta p = p_{0} - \frac{1}{3}p_{0} = \frac{2}{3}p_{0}$

根据动量守恒定律，甲、乙系统总动量守恒，乙与甲动量变化大小相等，乙的初动量为零，则刚分开时，乙的动量大小为 $\frac{2}{3}p_{0}$，所以甲、乙的动量大小之比为 1:2。

故选 C。
}

\item
%% number: 3
%% paperName: 2025年贵州高考第 3 题 4 分
%% typeId: 1
%% type: 单选题
%% chapter: 光学
%% point: 光的偏振
%% method: 无
%% score: 4
%% degree: 750
%% duplicateId: 0
%% body:
天花板上有一发光的白炽灯，桌上有一盆水。某同学第一次透过偏振片观察白炽灯，如图~\ref{2025贵州03a}所示；第二次透过偏振片观察白炽灯在水中的倒影，如图~\ref{2025贵州03b}所示。在两次观察中，以所观察光的传播方向为轴旋转偏振片时，透过偏振片观察到\xzanswer{C}

\twopicture[labelA=2025贵州03a, labelB=2025贵州03b, wA=2.7cm, wB=2.9cm]{figs/03a.png}{figs/03b.png}

\fourchoices[answer=C]
{白炽灯和倒影的亮度均变化}
{白炽灯和倒影的亮度均不变}
{白炽灯亮度不变，倒影亮度变化}
{白炽灯亮度变化，倒影亮度不变}

%% answer: C
%% memo:
\memoanswer{
白炽灯发出的光为自然光，其振动方向均匀分布，所以第一次透过偏振片直接观察白炽灯，白炽灯亮度不变；而自然光通过水面反射后为部分偏振光，则第二次透过偏振片观察白炽灯在水中的倒影时，当旋转偏振片时，偏振片的透振方向与反射光的主要振动方向的夹角不断变化，所以倒影的亮度会发生变化。

故选 C。
}

\item
%% number: 4
%% paperName: 2025年贵州高考第 4 题 4 分
%% typeId: 1
%% type: 单选题
%% chapter: 力物体的平衡
%% point: 三力平衡
%% method: 无
%% score: 4
%% degree: 650
%% duplicateId: 0
%% body:
一不可伸长的轻绳跨过同一水平线上的定滑轮，中间两定滑轮的间距为 $d$，在绳中央固定有一轻质吊环，绳两端分别挂有质量均为 $m$ 的配重物，配重物静止在地面上且绳恰好伸直。如图，某同学在健身时把吊环竖直向下缓慢拉 $\frac{d}{2}$ 的距离后保持静止，已知重力加速度大小为 $g$，不计摩擦及滑轮大小，则静止时该同学对吊环的拉力大小为\xzanswer{A}

\onepicture[w=4.2cm]{figs/04.png}

\fourchoices[answer=A]
{$\sqrt 2 mg$}
{$mg$}
{$\frac{\sqrt 2}{2}mg$}
{$\frac{1}{2}mg$}

%% answer: A
%% memo:
\memoanswer{
根据几何关系可知，吊环两边的细绳与竖直方向的夹角为 $45^\circ$，则由平衡可知

$F = 2T\cos 45^\circ$

其中 $T = mg$，可得 $F = \sqrt 2 mg$

故选 A。
}

\item
%% number: 5
%% paperName: 2025年贵州高考第 5 题 4 分
%% typeId: 1
%% type: 单选题
%% chapter: 电场
%% point: 电场线
%% method: 无
%% score: 4
%% degree: 650
%% duplicateId: 0
%% body:
如图，$M$ 和 $N$ 是在一直线上相距 $L$ 的两点，$A$ 是 $MN$ 的中点，直线外的 $B$ 点与 $N$ 点相距 $\frac{L}{2}$。在 $M$、$N$ 两点分别固定等量的正点电荷 $q$，选取无穷远处为零电势点，设 $A$、$B$ 两点的电场强度大小分别为 $E_{A}$、$E_{B}$，电势分别为 $\varphi_{A}$、$\varphi_{B}$，则\xzanswer{A}

\onepicture[w=5.0cm]{figs/05.png}

\fourchoices[answer=A]
{$E_{A} < E_{B}$，$\varphi_{A} > \varphi_{B}$}
{$E_{A} < E_{B}$，$\varphi_{A} < \varphi_{B}$}
{$E_{A} > E_{B}$，$\varphi_{A} > \varphi_{B}$}
{$E_{A} > E_{B}$，$\varphi_{A} < \varphi_{B}$}

%% answer: A
%% memo:
\memoanswer{
根据 $E = k \frac{Q}{r^{2}}$，结合矢量合成可知，两电荷在 $A$ 点产生场强等大反向，所以 $E_{A} = 0$。而两电荷在 $B$ 点合场强不为零，所以 $E_{A} < E_{B}$。等量同种点电荷的电场分布，如图所示。

\onepicture[w=3.6cm]{figs/05_an.png}

根据沿着电场线方向电势逐渐降低，可看出 $\varphi_{A} > \varphi_{B}$。

故选 A。
}

\item
%% number: 6
%% paperName: 2025年贵州高考第 6 题 4 分
%% typeId: 1
%% type: 单选题
%% chapter: 曲线运动
%% point: 人造地球卫星
%% method: 无
%% score: 4
%% degree: 550
%% duplicateId: 0
%% body:
一质量为 $m$ 的人造卫星绕地球做轨道半径为 $R$ 的匀速圆周运动。由于存在稀薄空气，经过一段时间后，卫星做圆周运动的轨道半径变为 $kR$。已知地球的质量为 $M$，引力常量为 $G$，则在该段时间内人造卫星所受力的合力做的功为\xzanswer{B}

\fourchoices[answer=B]
{$\frac{GMm(k - 1)}{2R}$}
{$\frac{GMm(1 - k)}{2kR}$}
{$\frac{GMm(k - 1)}{R}$}
{$\frac{GMm(1 - k)}{kR}$}

%% answer: B
%% memo:
\memoanswer{
根据动能定理，合外力做功等于物体动能的变化量，即 $W_{\text{合}} = \Delta E_{\text{k}}$

卫星绕地球做匀速圆周运动时，万有引力提供向心力，由 $G\frac{Mm}{r^{2}} = m\frac{v^{2}}{r}$

可得卫星动能 $E_{\text{k}} = \frac{1}{2}mv^{2} = \frac{GMm}{2r}$

初态轨道半径为 $R$，初动能 $E_{\text{k}1} = \frac{GMm}{2R}$

末态轨道半径为 $kR$，末动能 $E_{\text{k}2} = \frac{GMm}{2kR}$

代入动能定理计算得 $W_{\text{合}} = \frac{GMm}{2kR} - \frac{GMm}{2R} = \frac{GMm(1 - k)}{2kR}$

故选 B。
}

\item
%% number: 7
%% paperName: 2025年贵州高考第 7 题 4 分
%% typeId: 1
%% type: 单选题
%% chapter: 交流电
%% point: 交流电
%% method: 无
%% score: 4
%% degree: 450
%% duplicateId: 0
%% body:
某同学将一底面半径为 $r$、高为 $L$ 的圆柱形绝缘竹笼改造成仓鼠玩具。如图，在竹笼柱面上平行于轴线固定有一长为 $L$ 的轻质金属棒，其两端通过导线连接一电阻为 $R$ 的小灯泡，竹笼所在区域存在竖直向上、磁感应强度大小为 $B$ 的匀强磁场。若仓鼠在竹笼中奔跑时，竹笼绕水平中轴线以角速度 $\omega$ 匀速转动，小灯泡电阻视为不变，不计其他电阻，则竹笼转动一圈小灯泡消耗的电能为\xzanswer{B}

\onepicture[w=4.6cm]{figs/07.png}

\fourchoices[answer=B]
{$\frac{2\pi \omega B^{2}L^{2}r^{2}}{R}$}
{$\frac{\pi \omega B^{2}L^{2}r^{2}}{R}$}
{$\frac{\pi \omega B^{2}L^{4}}{2R}$}
{$\frac{\pi \omega B^{2}r^{4}}{2R}$}

%% answer: B
%% memo:
\memoanswer{
根据正弦式交流电产生原理可知，交流电最大值为 $E_{m} = BLr\omega$

电压有效值为 $U = \frac{E_{m}}{\sqrt{2}}$

竹笼转动一圈小灯泡消耗的电能为 $E = \frac{U^{2}}{R} \times \frac{2\pi}{\omega} = \frac{\pi \omega B^{2}L^{2}r^{2}}{R}$

故选 B。
}

\item
%% number: 8
%% paperName: 2025年贵州高考第 8 题 5 分
%% typeId: 2
%% type: 多选题
%% chapter: 光学
%% point: 光子说
%% method: 无
%% score: 5
%% degree: 650
%% duplicateId: 0
%% body:
观测分子云中的 CO 射电谱线对我们理解恒星的诞生、星际物质循环等具有重要意义。我国正在开展的第二期“银河画卷”巡天项目中，观测 \ce{^{12}C^{16}O}、\ce{^{13}C^{16}O}、\ce{^{12}C^{18}O} 同位素分子辐射出的电磁波的频率如表所示。下列说法正确的是\xzanswer{BD}

\begin{center}
\begin{tabular}{|c|c|}
\hline
CO同位素分子 & 频率/（$10^{11} \UHz$） \\
\hline
\ce{^{12}C^{16}O} & 1.153 \\
\hline
\ce{^{13}C^{16}O} & 1.102 \\
\hline
\ce{^{12}C^{18}O} & 1.098 \\
\hline
\end{tabular}
\end{center}

\fourchoices[answer=BD]
{\ce{^{12}C^{18}O} 比 \ce{^{12}C^{16}O} 多两个质子}
{\ce{^{13}C^{16}O} 比 \ce{^{12}C^{16}O} 多一个中子}
{\ce{^{12}C^{16}O} 辐射出的光子的能量最小}
{\ce{^{12}C^{16}O} 辐射出的电磁波的波长最短}

%% answer: BD
%% memo:
\memoanswer{
A．\ce{^{12}C^{18}O} 和 \ce{^{12}C^{16}O} 中 C 均为 \ce{^{12}C}，质子数均为 6，氧的质子数均为 8，两者分子总质子数均为 $6 + 8 = 14$，质子数相等，故 A 错误；

B．\ce{^{13}C} 中子数为 $13 - 6 = 7$，\ce{^{12}C} 中子数为 $12 - 6 = 6$，\ce{^{13}C} 比 \ce{^{12}C} 多 1 个中子，两者的 O 均为 \ce{^{16}O}，中子数相同，故 \ce{^{13}C^{16}O} 总中子数比 \ce{^{12}C^{16}O} 多 1，故 B 正确；

C．光子能量公式为 $E = h\nu$，由表格可知 \ce{^{12}C^{16}O} 的辐射频率最大，对应光子能量最大，故 C 错误；

D．电磁波波长与频率关系为 $\lambda = \frac{c}{\nu}$，频率越大波长越短，\ce{^{12}C^{16}O} 频率最大，故波长最短，故 D 正确。

故选 BD。
}

\item
%% number: 9
%% paperName: 2025年贵州高考第 9 题 5 分
%% typeId: 2
%% type: 多选题
%% chapter: 气体性质
%% point: 热力学第一定律
%% method: 无
%% score: 5
%% degree: 450
%% duplicateId: 0
%% body:
如图，在固定的绝热气缸中固定一多孔塞，其两侧各有一绝热活塞。开始时，多孔塞左侧有一定质量的理想气体，其状态参量为 $p_{1}$、$V_{1}$、$T_{1}$，内能为 $U_{1}$，右侧活塞紧靠多孔塞。在外力 $F_{1}$ 作用下，左侧气体在压强 $p_{1}$ 不变的条件下缓慢通过多孔塞流向右侧，右侧气体压强因外力 $F_{2}$ 作用始终为 $p_{2}$。当气体全部流入右侧后，其状态参量为 $p_{2}$、$V_{2}$、$T_{2}$，内能为 $U_{2}$。已知 $p_{1} > p_{2}$，$\frac{U_{1}}{T_{1}} = \frac{U_{2}}{T_{2}}$，不计活塞与气缸间的摩擦，装置密闭性良好，则\xzanswer{AC}

\onepicture[w=4.0cm]{figs/09.png}

\fourchoices[answer=AC]
{整个过程中左侧活塞对气缸中气体做功为 $p_{1}V_{1}$}
{整个过程中气缸中气体对右侧活塞做功为 $p_{1}V_{2}$}
{$U_{1} + p_{1}V_{1} = U_{2} + p_{2}V_{2}$}
{$T_{1} < T_{2}$}

%% answer: AC
%% memo:
\memoanswer{
A．气体缓慢通过多孔塞，设圆柱的横截面积为 $S$，所以整个过程中左侧活塞对气缸中气体做功为 $W = F_{1}L_{1} = p_{1}SL_{1} = p_{1}V_{1}$，故 A 正确；

B．同理，整个过程中气缸中气体对右侧活塞做功为 $W' = p_{2}\Delta V = p_{2}V_{2}$，故 B 错误；

C．固定的绝热气缸两侧有绝热活塞，该系统为绝热系统即 $Q = 0$，根据热力学第一定律可知 $\Delta U = W$

即 $U_{2} - U_{1} = p_{1}V_{1} + (-p_{2}V_{2})$

整理得 $U_{1} + p_{1}V_{1} = U_{2} + p_{2}V_{2}$，故 C 正确；

D．内能与温度成正比 $U_{1} = kT_{1}$，$U_{2} = kT_{2}$

根据 $\frac{pV}{T} = C$ 得 $p_{1}V_{1} = CT_{1}$，$p_{2}V_{2} = CT_{2}$

代入方程 $U_{1} + p_{1}V_{1} = U_{2} + p_{2}V_{2}$

得 $kT_{1} + CT_{1} = kT_{2} + CT_{2}$

两边约掉 $(k + C)$ 得 $T_{1} = T_{2}$，故 D 错误。

故选 AC。
}

\item
%% number: 10
%% paperName: 2025年贵州高考第 10 题 5 分
%% typeId: 2
%% type: 多选题
%% chapter: 电磁感应
%% point: 电磁感应能量分析
%% method: 无
%% score: 5
%% degree: 450
%% duplicateId: 0
%% body:
如图~\ref{2025贵州10a}，一竖直固定的透明塑料管内固定有 6 个小磁铁，相邻磁铁同极靠近、间距很小，取距离最上端的小磁铁 S 极 $0.15 \Um$ 处为坐标原点 $O$，$y$ 轴正方向竖直向下。将一内径略大于塑料管外径的金属环套在塑料管上，在 $O$ 点处由静止释放，金属环的速度一位置图像如图~\ref{2025贵州10b}所示。已知金属环的质量为 $35.9 \Ug$，取重力加速度大小为 $10 \Umsq$，不计摩擦和空气阻力，则下落过程中，金属环\xzanswer{AC}

\twopicture[labelA=2025贵州10a, labelB=2025贵州10b, wA=1.9cm, wB=9.0cm]{figs/10a.png}{figs/10b.png}

\fourchoices[answer=AC]
{在 $0.11 \leq y \leq 0.15 \Um$ 内所受安培力方向竖直向上}
{在 $0.35 \leq y \leq 0.39 \Um$ 内所受安培力方向竖直向下}
{在 $0.24 \leq y \leq 0.26 \Um$ 内克服安培力做的功为 $7.1\times10^{-3} \UJ$}
{在 $0.11 \leq y \leq 0.15 \Um$ 内感应电流方向为顺时针方向（从上向下看）}

%% answer: AC
%% memo:
\memoanswer{
AD．在 $0.11 \leq y \leq 0.15 \Um$ 内，由图（a）可知，穿过金属环的磁通量向下增大，根据楞次定律可知，感应电流产生的磁场方向向上，由安培定则可知，感应电流方向为逆时针方向（从上向下看）；根据楞次定律“来拒去留”推论可知，金属环所受安培力方向竖直向上，故 A 正确，D 错误；

B．在 $0.35 \leq y \leq 0.39 \Um$ 内，穿过金属环的磁通量向上减小，根据楞次定律结合安培定则可知，感应电流方向为逆时针（从上向下看），根据楞次定律“来拒去留”推论可知，金属环所受安培力方向竖直向上，故 B 错误；

C．由图（b）可知，在 $0.24 \leq y \leq 0.26 \Um$ 内金属环的动能变化为 0，根据动能定理可得 $mgh - W = 0$，可得克服安培力做的功为 $W = mgh = 35.9 \times 10^{-3} \times 10 \times (0.26 - 0.24) \UJ = 7.18 \times 10^{-3} \UJ$，故 C 正确。

故选 AC。
}

\item
%% number: 11
%% paperName: 2025年贵州高考第 11 题 6 分
%% typeId: 4
%% type: 实验题
%% chapter: 曲线运动
%% point: 研究平抛运动实验
%% method: 无
%% score: 6
%% degree: 650
%% duplicateId: 0
%% body:
在“探究平抛运动的特点”实验中，某同学使用视频分析软件提取小球运动的位置数据。所用器材有高速摄像机、小钢球、斜槽轨道、铅垂线、刻度尺等。

请完成下列步骤：

\begin{enumerate}
\item
固定斜槽轨道并使其末端 \tkanswer{水平}。

\item
在斜槽轨道上某位置释放小钢球，用摄像机录制小钢球的运动过程。

\item
通过视频分析，提取一段时间内每隔 $0.05 \Us$ 小钢球的位置，测量得到小钢球相邻两个位置的水平间距 $\Delta x$ 和竖直间距 $\Delta y$ 如下表所示：

\begin{center}
\begin{tabular}{|c|c|c|c|c|}
\hline
$\Delta x/\Um$ & 0.0517 & 0.0514 & 0.0505 & 0.0504 \\
\hline
$\Delta y/\Um$ & 0.0975 & 0.1217 & 0.1465 & 0.1704 \\
\hline
\end{tabular}
\end{center}

经初步分析，可判断小钢球在水平方向做匀速直线运动、竖直方向做匀变速直线运动。充分利用数据计算得小钢球的水平方向速度大小为 \tkanswer{1.02} $\Ums$、竖直方向加速度大小为 \tkanswer{9.77} $\Umsq$。（结果均保留 2 位小数）
\end{enumerate}

%% answer:
\jdanswer{
\begin{enumerate}
\item
水平

\item
（本步骤无作答内容）

\item
1.02（1.01 或 1.03），9.77
\end{enumerate}
}
%% memo:
\memoanswer{
（1）固定斜槽轨道并使其末端水平，以保证小球做平抛运动；

（3）由表中数据可知，水平位移的平均值为
$x = \frac{0.0517 + 0.0514 + 0.0505 + 0.0504}{4} \Um = 0.051 \Um$

小钢球的水平方向速度大小 $v_{0} = \frac{x}{T} = \frac{0.051}{0.05} \Ums = 1.02 \Ums$

竖直方向根据 $\Delta y' = gT^{2}$，结合逐差法可得

$g = \frac{(\Delta y_{4} + \Delta y_{3}) - (\Delta y_{2} + \Delta y_{1})}{4T^{2}} = \frac{(0.1704 + 0.1465) - (0.1217 + 0.0975)}{4 \times 0.05^{2}} \Umsq = 9.77 \Umsq$
}

\item
%% number: 12
%% paperName: 2025年贵州高考第 12 题 9 分
%% typeId: 4
%% type: 实验题
%% chapter: 电路
%% point: 电路实验
%% method: 无
%% score: 9
%% degree: 450
%% duplicateId: 0
%% body:
某实验小组为测量粗细均匀的铅笔芯单位长度电阻及电压表的内阻，选用的器材有：

待测铅笔芯；

电源 $E$（电动势 $1.5 \UV$，内阻不计）；

毫安表 A（内阻约 $10 \UO$，量程 $0 \sim 100 \UmA$）；

电压表 V（内阻约几百欧，量程 $0 \sim 1 \UV$）；

滑动变阻器（阻值 $0 \sim 20 \UO$）；

刻度尺；开关；导线若干。

请完成下列步骤：

\begin{enumerate}
\item
该小组设计了如图~\ref{2025贵州12a}所示的电路图。根据图~\ref{2025贵州12a}在答题卡上完成图~\ref{2025贵州12b}中的实物图连线。

\item
闭合开关 S 前，应将滑动变阻器的滑片调到 \tkanswer{B} 端（填“A”或“B”）。

\item
移动滑动触头 P 到铅笔芯上的某位置，测量铅笔芯接入电路部分的长度 $x$；闭合 S，移动滑动变阻器的滑片使毫安表示数 $I_{0} = 50.0 \UmA$，读出此时电压表的示数 $U$，断开 S。

\item
重复步骤（3），得到多组 $I_{0} = 50.0 \UmA$ 条件下不同 $x$ 对应的 $U$，如下表所示。表中电压的某个数据记录有误，有误的数据是 \tkanswer{0.4}。

\begin{center}
\begin{tabular}{|c|c|c|c|c|c|}
\hline
$x/\Ucm$ & 3.00 & 4.00 & 6.00 & 10.00 & 14.00 \\
\hline
$U/\UV$ & 0.13 & 0.18 & 0.26 & 0.4 & 0.57 \\
\hline
\end{tabular}
\end{center}

\item
实验小组更正表中的数据后，计算出 $\frac{1}{x}$ 和 $\frac{1}{U}$ 的值，并以 $\frac{1}{x}$ 为横坐标、$\frac{1}{U}$ 为纵坐标对其进行线性拟合，得到直线的斜率 $k = 0.23 \Um\cdot\UV^{-1}$，纵轴截距 $b = 0.06 \UV^{-1}$。计算出该铅笔芯单位长度（$1 \Um$）的电阻为 \tkanswer{87} $\UO$，电压表的内阻为 \tkanswer{333} $\UO$。（结果均保留至整数）
\end{enumerate}

\twopicture[labelA=2025贵州12a, labelB=2025贵州12b, wA=4.0cm, wB=6.0cm, align=right]{figs/12a.png}{figs/12b.png}

%% answer:
\jdanswer{
\begin{enumerate}
\item
（连线图见详解）

\item
B

\item
（本步骤无作答内容）

\item
0.4

\item
87，333
\end{enumerate}
}
%% memo:
\memoanswer{
（1）根据电路图描绘实物图如下

\onepicture[w=5.0cm]{figs/12_an.png}

（2）滑动变阻器限流式连接，闭合开关 S 前，应将滑动变阻器的滑片调到 B 端使回路电阻最大。

（4）表格中电压表的读数均为保留两位小数，电压有误的数据是 $0.4 \UV$，应为 $0.40 \UV$。

（5）设该铅笔芯单位长度（$1 \Um$）的电阻为 $R_{0}$，根据电路图可知

$U = (I_{0} - \frac{U}{R_{\mathrm{V}}})R_{x} = (I_{0} - \frac{U}{R_{\mathrm{V}}}) \times xR_{0}$

整理得 $\frac{1}{U} = \frac{1}{x} \cdot \frac{1}{I_{0}R_{0}} + \frac{1}{I_{0}R_{\mathrm{V}}}$

已知 $k = \frac{1}{I_{0}R_{0}} = 0.23 \Um\cdot\UV^{-1}$，$b = \frac{1}{I_{0}R_{\mathrm{V}}} = 0.06 \UV^{-1}$

计算得 $R_{0} = 87 \UO$，$R_{\mathrm{V}} = 333 \UO$
}

\item
%% number: 13
%% paperName: 2025年贵州高考第 13 题 8 分
%% typeId: 6
%% type: 计算题
%% chapter: 机械波
%% point: 波与振动图象
%% method: 无
%% score: 8
%% degree: 550
%% duplicateId: 0
%% body:
贵州黎平的肇兴侗寨有“鼓楼之乡”的美誉，图~\ref{2025贵州13a}的建筑群展示了精湛的传统鼓楼木构技艺。为保护鼓楼等传统木质建筑，可采用超声横波法进行无损检测，获得其弹性模量等力学参数，从而采取相应的保护措施。在某次检测中，有一列简谐横波在木材中沿 $x$ 轴传播，$t = 0$ 时刻的波形图如图~\ref{2025贵州13b}所示，位于 $x$ 轴上的质点 $P$ 的振动图像如图~\ref{2025贵州13c}所示。求：

\threepicture[labelA=2025贵州13a, labelB=2025贵州13b, labelC=2025贵州13c, wA=3.2cm, wB=4.4cm, wC=3.4cm]{figs/13a.png}{figs/13b.png}{figs/13c.png}

\begin{enumerate}
\item
波在木材中的传播方向及速度；
\item
在 $0 \sim 3\times10^{-6} \Us$ 时间内，质点 $P$ 运动的路程。
\end{enumerate}

%% answer:
\jdanswer{
\begin{enumerate}
\item
传播方向沿 $x$ 轴负方向，$800 \Ums$

\item
$1.2\times10^{-5} \Um$
\end{enumerate}
}
%% memo:
\memoanswer{
（1）由 P 点的振动图像可知，$t = 0$ 时刻质点 $P$ 在平衡位置沿 $y$ 轴负向振动，结合波形图可知，波传播方向沿 $x$ 轴负方向；

根据波形图可知波长 $\lambda = 8\times10^{-4} \Um$，根据振动图像可知周期 $T = 1\times10^{-6} \Us$。

波在木材中的传播速度 $v = \frac{\lambda}{T} = 800 \Ums$

（2）时间关系 $3\times10^{-6} \Us = 3T$，可知 $0 \sim 3\times10^{-6} \Us$ 时间内，质点 $P$ 完成了 3 个全振动，故其运动的路程 $s = 3\times4A = 12\times1\times10^{-6} \Um = 1.2\times10^{-5} \Um$
}

\item
%% number: 14
%% paperName: 2025年贵州高考第 14 题 15 分
%% typeId: 6
%% type: 计算题
%% chapter: 运动和力
%% point: 牛顿定律的应用
%% method: 无
%% score: 15
%% degree: 450
%% duplicateId: 0
%% body:
杵臼是我国古代加工谷物的重要工具，在《诗经$\cdot$大雅》中有明确记载。使用杵臼的示意图如图~\ref{2025贵州14a}所示，舂捣臼中谷物时，手紧握质量为 $m$ 的石杵（石杵与谷物接触但未陷入），对其施加一竖直向上的恒力 $F_{1}$ 使其上升，作用一段时间 $t_{1}$ 后松手，松手后不考虑手与石杵的作用力。当石杵上升到最高点时，手再次紧握石杵并对其施加一竖直向下的作用力 $F_{2}$，其大小随下降距离 $x$ 的变化关系如图~\ref{2025贵州14b}所示，图中 $F_{m}$ 为 $F_{2}$ 的最大值。石杵接触谷物时松手，松手后不考虑手与石杵的作用力，再经过 $\Delta t$ 时间石杵静止，完成一次舂捣。已知 $m = 2.5 \Ukg$，$F_{1} = \frac{200}{7} \UN$，$F_{m} = 50 \UN$，$t_{1} = 0.7 \Us$，$\Delta t = 0.025 \Us$，取重力加速度大小 $g = 10 \Umsq$。求：

\begin{enumerate}
\item
石杵上升的最大高度 $h_{0}$ 及上升过程所用的时间；
\item
$\Delta t$ 时间内石杵对谷物的平均作用力大小。
\end{enumerate}

\twopicture[labelA=2025贵州14a, labelB=2025贵州14b, wA=2.7cm, wB=3.6cm, align=right]{figs/14a.png}{figs/14b.png}

%% answer:
\jdanswer{
\begin{enumerate}
\item
$0.4 \Um$，$0.8 \Us$

\item
$425 \UN$
\end{enumerate}
}
%% memo:
\memoanswer{
（1）对石杵施加一竖直向上的恒力 $F_{1}$，当作用时间为 $t_{1}$ 的过程中的加速度

$a = \frac{F_{1} - mg}{m} = \frac{10}{7} \Umsq$

此时的速度 $v_{1} = at_{1} = 1 \Ums$

上升的位移 $h_{1} = \frac{v_{1}}{2}t_{1} = 0.35 \Um$

撤去 $F_{1}$ 后还能上升的高度 $h_{2} = \frac{v_{1}^{2}}{2g} = 0.05 \Um$

还能上升的时间 $t_{2} = \frac{v_{1}}{g} = 0.1 \Us$

石杵上升的最大高度 $h_{0} = h_{1} + h_{2} = 0.4 \Um$

上升过程所用的时间 $t_{0} = t_{1} + t_{2} = 0.8 \Us$

（2）根据图像，石杵下落过程中 $F_{2}$ 对石杵做功为

$W = \frac{1}{2}F_{m}h_{0} = \frac{1}{2} \times 50 \times 0.4 \UJ = 10 \UJ$

当到达石杵接触谷物时由动能定理

$W + mgh_{0} = \frac{1}{2}mv_{2}^{2}$

解得 $v_{2} = 4 \Ums$

石杵与谷物作用的过程，对石杵由动量定理（向上为正）$(F - mg)\Delta t = 0 - (-mv_{2})$

解得 $F = 425 \UN$

根据牛顿第三定律可知，石杵对谷物的平均作用力大小 $425 \UN$。
}

\item
%% number: 15
%% paperName: 2025年贵州高考第 15 题 19 分
%% typeId: 6
%% type: 计算题
%% chapter: 磁场
%% point: 带电粒子在复合场中的运动
%% method: 无
%% score: 19
%% degree: 350
%% duplicateId: 0
%% body:
如图，建立直角坐标系 $O-xyz$，$x$ 轴正方向水平向右，$y$ 轴正方向垂直纸面向里（$y$ 轴未画出），$z$ 轴正方向竖直向上。空间中存在方向竖直向上的匀强电场，在 $z \leq z_{0}$ 的区域 Ⅰ 和 $z > z_{0}$ 的区域 Ⅱ 中均存在方向垂直纸面向里的匀强磁场，区域 Ⅰ 的磁感应强度大小为 $B_{1}$，区域 Ⅱ 的磁感应强度大小 $B_{2}$ 未知。有一带正电荷的粒子，质量为 $m$、电荷量为 $q$，以速率 $v$ 从 $O$ 点沿 $x$ 轴正方向射出后，在区域 Ⅰ、Ⅱ 中均可做匀速圆周运动，且恰好能经过 $x$ 轴上的 P 点，P 点坐标为（$\frac{16v^{2}}{5g}$，0，0）。已知 $z_{0} = \frac{2v^{2}}{5g}$，$B_{1} = \frac{mg}{qv}$，$g$ 为重力加速度。

\begin{enumerate}
\item
求电场强度大小 $E$ 及该粒子第一次经过 $z = z_{0}$ 平面时的位置对应的 $x$ 坐标值；
\item
求粒子从 $O$ 点到 P 点的运动时间最短时区域 Ⅱ 的磁感应强度大小 $B_{2}$；
\item
若仅将匀强电场的方向改为沿 $y$ 轴正向，该粒子仍以速率 $v$ 从 $O$ 点沿 $x$ 轴正方向射出，求该粒子的轨迹方程。
\end{enumerate}

\onepicture[align=right, w=6.5cm]{figs/15.png}

%% answer:
\jdanswer{
\begin{enumerate}
\item
$E = \frac{mg}{q}$，$x = \frac{4v^{2}}{5g}$

\item
$B_{2} = \frac{3mg}{qv}$

\item
$y = \frac{gx^{2}}{2v^{2}}$
\end{enumerate}
}
%% memo:
\memoanswer{
（1）由题意可知，粒子受到重力、洛伦兹力和电场力做匀速圆周运动，可以判断粒子受到的电场力与重力平衡，则有 $Eq = mg$

解得 $E = \frac{mg}{q}$

粒子在 Ⅰ 区做匀速圆周运动，运动轨迹如图所示

\onepicture[w=4.0cm]{figs/15_an1.png}

根据洛伦兹力提供向心力有 $qvB_{1} = m\frac{v^{2}}{R_{1}}$

解得 $R_{1} = \frac{v^{2}}{g}$

又根据几何关系有 $x^{2} + (R_{1} - z_{0})^{2} = R_{1}^{2}$

解得 $x = \frac{4v^{2}}{5g}$

（2）粒子做匀速圆周运动，可能的运动轨迹如图所示

\onepicture[w=6.0cm]{figs/15_an2.png}

设粒子进入 $B_{2}$ 磁场中速度方向与磁场分界面成 $\theta$ 角，根据几何关系有 $\cos(90^\circ - \theta) = \frac{x}{R_{1}}$

解得 $\theta = 53^\circ$

设粒子在 $B_{2}$ 磁场中运动的轨道半径为 $R_{2}$，根据圆周运动轨迹可知粒子运动到 P 点应满足

$n(2R_{1}\sin\theta - 2R_{2}\sin\theta) = \frac{16v^{2}}{5g}$

当 $n$ 取最小值时，运动时间最短。结合上图分析，可知带电粒子在 $B_{2}$ 磁场中至少绕 3 次才能到达 P 点，环绕的次数越少，用时越短，即 $n = 3$ 时所用的时间最短，则有

$3 \times (2R_{1}\sin\theta - 2R_{2}\sin\theta) = \frac{16v^{2}}{5g}$

解得 $R_{2} = \frac{v^{2}}{3g}$

根据洛伦兹力提供向心力有 $qvB_{2} = m\frac{v^{2}}{R_{2}}$

解得 $B_{2} = \frac{3mg}{qv}$

（3）若将电场方向改为 $y$ 轴方向正方向，由受力分析，粒子受到沿 $z$ 轴正方向的洛伦兹力、沿 $z$ 轴负方向的重力、沿 $y$ 轴正方向的电场力，则粒子在 Ⅰ 区受到的洛伦兹力大小为

$f = qvB_{1} = mg$

正好与重力相平衡，所以粒子在 $y$ 轴正方向做匀加速直线运动，有 $y = \frac{1}{2}at^{2}$

根据牛顿第二定律有 $Eq = ma$

解得 $a = g$

粒子在 $x$ 轴正方向做匀速直线运动，有 $x = vt$

联立解得 $y = \frac{gx^{2}}{2v^{2}}$
}

\end{enumerate}

\end{document}
